\documentclass[11pt]{article}

\usepackage[margin=1in]{geometry}
\usepackage{graphicx}
\usepackage{booktabs}
\usepackage{hyperref}
\usepackage{amsmath}
\usepackage{caption}
\usepackage{subcaption}
\usepackage{microtype}
\usepackage{authblk}

\title{Plant Geopriors: Late-Fusion of Visual and Geographic-Temporal Context for Fine-Grained Plant Species Classification}

\author[1]{Neha Sahha\thanks{author.one@placeholder.edu}}
\author[1]{Cecilia Zhang\}
\author[1]{Manini Banerjee\}
\affil[1]{Department of Data Science, Placeholder University}

\date{May 2026}

\begin{document}

\maketitle

\begin{abstract}
Fine-grained plant species classification from photographs is hard because many species are visually near-identical yet ecologically distinct. We test whether geographic and seasonal priors --- latitude, longitude, and day-of-year --- improve classification when fused with a CNN trained on photographs. Working with 44 species across three California families (Asteraceae, Poaceae, Apiaceae) drawn from GBIF / iNaturalist 2021, we train two single-modality baselines (a fine-tuned ResNet50 on 4{,}400 image-backed observations and a small MLP on 103{,}314 sinusoidally-encoded geo/time observations) and two late-fusion models that operate on baseline logits: a closed-form weighted addition $\alpha \cdot \text{image} + \beta \cdot \text{context}$ tuned on the validation split, and a trainable fusion MLP head over concatenated logits. On the image-backed test split (660 examples), the weighted-addition model reaches 47.9\% accuracy versus 45.9\% for the image-only baseline, and the trainable fusion MLP reaches 48.6\% balanced accuracy. Per-family analysis shows context-driven gains concentrate in Apiaceae (small, geographically structured species) while several widespread Poaceae species regress under fusion. Results suggest geographic priors meaningfully help on niche-restricted taxa but offer limited benefit when species are cosmopolitan; we discuss implications for richer environmental features and joint training.
\end{abstract}

\section{Background and Motivation}

Citizen-science platforms such as iNaturalist now produce tens of millions of geotagged plant observations per year, fueling computer-vision benchmarks like iNaturalist 2021 and field-deployed identification apps. Modern CNNs achieve strong top-1 accuracy on coarse taxonomy, but performance degrades sharply for fine-grained discrimination between sister species that share floral and vegetative morphology. Within California flora, three families illustrate this problem acutely: Asteraceae contains many small yellow composites that are difficult to separate visually even for trained botanists; Poaceae (grasses) are notoriously hard from photographs; and Apiaceae includes several umbel-bearing species with near-identical inflorescences.

Photographs are not the only signal in an iNaturalist record. Each observation also carries GPS coordinates and a timestamp. Many species occupy distinct ecological niches or have well-separated phenologies, so the joint distribution of latitude, longitude, and day-of-year carries discriminative information that the photograph alone may not. The question is whether this metadata, treated as a prior, provides a useful complement to image-based classification - and where it benefits the most.

\section{Problem Statement}

We frame the task as 44-way single-label classification. Given an observation $x = (\text{image}, \text{lat}, \text{lon}, \text{doy})$, predict the species $y \in \mathcal{Y}$ where $|\mathcal{Y}| = 44$. We compare four models: an image-only ResNet50, a context-only MLP over six sinusoidal features, a weighted-addition late fusion of the two baselines, and a trainable fusion MLP head over concatenated logits. The primary metric is accuracy on a held-out test split; we also report per-species and per-family accuracy to surface where each modality helps or hurts.

\section{Data and Exploratory Analysis}

\paragraph{Source.} Observations are drawn from a California-restricted GBIF / iNaturalist 2021 export covering Asteraceae, Poaceae, and Apiaceae. After filtering to species with at least 300 observations and sufficient seasonal and geographic spread, 44 species are retained.

\paragraph{Species selection.} Notebook \texttt{04\_species\_comparison.ipynb} compares three procedural shortlisting strategies --- raw count (\texttt{top\_n}), seasonal coverage (\texttt{time\_diverse}), and joint temporal-spatial diversity via $k$-means clustering on (lat, lon, doy) (\texttt{time\_geo\_diverse}). The third strategy was adopted because it spreads training mass across the species range rather than concentrating it on a few hotspot counties. The notebook also documents the ``yellow-flower problem'': roughly half of the Asteraceae shortlist are small yellow composites that are visually near-identical, motivating the search for non-visual signal.

\paragraph{Two working datasets.} The project maintains two aligned datasets with the same 44-class label space.

\begin{itemize}
    \item The \textbf{image-backed} dataset (\texttt{master\_observations.csv}, 4{,}400 rows; \texttt{split\_iid.csv} 70/15/15) contains every observation for which a usable iNaturalist image was downloaded. Splits are 3{,}080 train / 660 val / 660 test.
    \item The \textbf{expanded} dataset (\texttt{master\_observations\_geo\_expanded.csv}, 103{,}314 rows; \texttt{split\_geo\_expanded.csv} 70/15/15) is built by \texttt{scripts/build\_expanded\_splits.py}. It preserves the original 4{,}400-row split assignments verbatim (no leakage) and stratifies the 98{,}914 newly-added image-less observations by species with seed 42. Split sizes are 72{,}319 / 15{,}497 / 15{,}498. A manifest records SHA256 hashes and per-split row counts for reproducibility.
\end{itemize}

\paragraph{Feature engineering.} Six geographic-temporal features are derived from the raw fields:
\[
\begin{aligned}
&\text{lat}_{\sin/\cos} = \sin/\cos\!\left(\tfrac{\pi \cdot \text{lat}}{90}\right), \quad
\text{lon}_{\sin/\cos} = \sin/\cos\!\left(\tfrac{\pi \cdot \text{lon}}{180}\right),\\
&\text{doy}_{\sin/\cos} = \sin/\cos\!\left(\tfrac{\pi \cdot (\text{doy}-183)}{183}\right).
\end{aligned}
\]
This is intentionally minimal: no bioclim, no elevation, no distance-to-feature. The sinusoidal encoding preserves periodicity for day-of-year and provides smooth, bounded inputs to the MLP.

\paragraph{EDA observations.} Observation density peaks in late spring and early summer (DOY 100--250), reflecting both flowering phenology and citizen-scientist activity. Per-species counts are highly imbalanced --- Baccharis pilularis dominates with $\sim$1{,}400 observations while several Apiaceae species fall under 200. We mitigate imbalance during training with weighted cross-entropy and per-epoch per-species capping rather than oversampling.



\section{Methods and Models}

\paragraph{Image-only baseline.} A torchvision ResNet50 initialized with ImageNet weights, frozen except for a replaced final fully-connected layer of dimension $2048 \rightarrow 44$. Training uses standard ImageNet augmentations, weighted cross-entropy on the inverse capped frequencies, and Adam at lr$=$1e-3. The best validation accuracy reached during fine-tuning was 0.448; test-set accuracy is 0.459 (notebook 07, IID test split).

\paragraph{Context-only baseline.} A four-layer MLP with three hidden blocks of width 256, ReLU activations and dropout 0.3, mapping the six geo/time features to 44 logits. Trained twice: once on the 4{,}400-row image-backed split (test accuracy $\sim$0.07, only $\sim$3$\times$ random) and once on the expanded 103k-row split (validation accuracy 0.138). The expanded model yields strikingly uneven per-species behavior: Sanicula arctopoides reaches 67\% top-1, while 21 species sit at 0\% --- a signature of taxa whose ranges overlap completely.

\paragraph{Weighted-addition fusion (notebook 07).} Both baselines are loaded frozen. We compute the linear combination
\[
\text{final\_logits} = \alpha \cdot \text{image\_logits} + \beta \cdot \text{context\_logits}
\]
on the validation split, sweeping $\beta \in [0,1]$ on a 21-point grid with $\alpha = 1 - \beta$. The chosen $(\alpha, \beta)$ is then evaluated once on the test split. This model has zero trainable parameters; it is essentially a calibration of how much the geographic prior should shift the visual posterior. The selected weights for the v1 reported run favor the image side; \textit{[placeholder: insert exact $\alpha, \beta$ from the run --- visible in the validation grid table in notebook 07]}.

\paragraph{Trainable fusion MLP (notebook 08).} A single hidden-layer head over concatenated baseline logits:
\[
\mathrm{FusionMLP}: \; \mathbb{R}^{88} \xrightarrow{\text{Linear}} \mathbb{R}^{256} \xrightarrow{\text{ReLU}, \text{Dropout}} \mathbb{R}^{44}.
\]
Both backbones are frozen; only the head trains. Training operates on pre-cached logit tensors over the expanded geo dataset's image-having rows, with batch size 256, Adam at lr$=$1e-3, per-epoch per-species capping at 500, and early stopping on validation balanced accuracy (best epoch 17 of a 30-epoch budget; stopped at epoch 22). Because backbones are frozen and the input is 88-d, the head is small ($\sim$25k parameters) and trains in seconds.

\paragraph{Reproducibility.} A single canonical pipeline script \texttt{scripts/build\_expanded\_splits.py} regenerates splits and feature tables from \texttt{time\_geo\_species\_df.csv} with seed 42 and emits a manifest of SHA256 hashes. Notebook 07 saves \texttt{weighted\_addition\_logits.pth} (val and test logits, labels, gbifIDs, chosen weights) so notebook 08 can rebuild every comparison plot without re-running the ResNet50.

\section{Results and Interpretation}

\paragraph{Overall test accuracy.} Table~\ref{tab:overall} summarizes top-line numbers. Both fusion models beat the image-only baseline; gains are small but consistent.

\begin{table}[h]
\centering
\caption{Test-set accuracy on the 660-row image-backed IID test split. ``Balanced'' denotes per-class recall averaged over species.}
\label{tab:overall}
\begin{tabular}{lcc}
\toprule
Model & Accuracy & Loss \\
\midrule
Context only (Geo MLP) & 0.056 & 3.63 \\
Image only (ResNet50)  & 0.459 & 2.11 \\
Weighted addition (07) & 0.479 & 2.38 \\
Fusion MLP (08, balanced) & 0.486 & --- \\
\bottomrule
\end{tabular}
\end{table}

The image baseline carries most of the signal: context alone is barely better than chance for the 4{,}400-row geo model and reaches 13.8\% validation accuracy on the expanded data. Both fusion approaches improve over image-only by 2--3 percentage points, with the trainable head edging the closed-form weighted sum.

\paragraph{Per-species structure.} Figure~\texttt{fusion\_mlp\_per\_species\_accuracy\_TESTSET.png} and Figure~\texttt{accuracy\_delta\_TESTSET.png} (in \texttt{results/}) show that the headline averages mask substantial per-species variance. The fusion MLP exceeds 80\% accuracy on Sanicula arctopoides (87\%) and Nassella pulchra (80\%) while Achillea millefolium falls to 7\%. Twenty-one of forty-four species clear 50\% accuracy under fusion, and no species sits at exactly zero --- an improvement over the geo-only model.

\paragraph{Where fusion helps and where it hurts.} A regression analysis in notebook 08 isolates seven species where the fusion MLP underperforms the ResNet50 baseline: \textit{Arundo donax}, \textit{Encelia farinosa}, \textit{Distichlis spicata}, \textit{Centaurea solstitialis}, \textit{Briza maxima}, \textit{Cynosurus echinatus}, and \textit{Bromus hordeaceus}. These are mostly widespread weedy or grass species whose ranges blanket the state. For such taxa the geographic prior is essentially uninformative and adds noise that occasionally pulls the argmax away from a confident image prediction. Conversely, the most-improved species (visualized in \texttt{most\_improved\_species.png}) are species with restricted, well-separated ranges where the geographic prior breaks ties between visually similar look-alikes.

\paragraph{Family breakdown.} \texttt{family\_breakdown\_TEST.png} reports balanced accuracy grouped by family. Apiaceae is the easiest family across all four models (fusion MLP $\sim$0.87), reflecting both small species count within the family and strong geographic separation between umbellifers. Asteraceae sits in the middle and Poaceae is the hardest --- consistent with the notebook 04 finding that grasses are both visually impoverished from photographs and geographically generalist.

\paragraph{Confusions.} \texttt{confused\_pairs\_fusion\_mlp.png} surfaces the top confusions remaining after fusion. The pattern is dominated by within-genus pairs (e.g.\ \textit{Sanicula} cohabitants, congeneric \textit{Bromus} grasses) where image, geography, and phenology are all close --- a regime where neither modality has a clean advantage and richer features (within-image attention, bioclim variables, microhabitat) would likely be required.

\paragraph{Confusion matrices.} \texttt{all\_models\_confusion\_matrices\_TEST.png} places the four models side-by-side. The image-only and fusion matrices share dense diagonals; the geo-only matrix shows a few strong diagonal blocks (range-restricted species) embedded in heavy off-diagonal noise. Weighted addition and the fusion MLP produce visibly cleaner matrices than the image-only baseline, with the fusion MLP slightly tighter.

\paragraph{Caveats on the comparison.} Notebooks 07 and 08 evaluate on the same 660-row image-backed test set, but use slightly different metric definitions: notebook 07 reports micro accuracy, notebook 08 reports balanced (macro-recall) accuracy. The numbers in Table~\ref{tab:overall} preserve the original definitions; like-for-like comparison should recompute both from the saved logits. A small calibration analysis (expected calibration error, reliability diagrams) would also be informative but is \textit{[placeholder: not yet completed]}.

\section{Conclusions and Discussion}

Geographic and seasonal context, even reduced to six sinusoidal features, provides a measurable but modest boost over a strong image baseline on this 44-species California subset. Two simple late-fusion designs --- closed-form weighted addition and a small trained head over logits --- both outperform the image-only model, with the trained head marginally ahead. The benefit concentrates in species with restricted ranges and clean phenology, while widespread weedy species regress under fusion because the geographic prior is uninformative for them.

Three follow-ups stand out. First, the feature set is unusually minimal; substituting or augmenting the six sinusoidal features with bioclim variables, elevation, or species-distribution-model priors would test whether the gains scale with feature richness. Second, both fusion models keep the backbones frozen; an end-to-end fine-tune of the ResNet50 with concatenated context features would let the visual representation adapt to the prior. Third, calibration is unmeasured: even where accuracy is unchanged, fusion may improve confidence calibration, which matters in deployed identification tools where a confident-but-wrong prediction is worse than a hedged one.

\section{Broader Impact}

Plant identification tools shape how citizen scientists, land managers, and educators perceive biodiversity. False positives on widespread weedy species risk inflating reports of invasive presence; false negatives on rare natives risk underestimating conservation value. Fusing geographic priors helps the easy cases (range-restricted natives) but can mask the hard cases (cosmopolitan invasives) by collapsing them onto more common neighbors. Any deployment should expose calibration to the user and refrain from over-claiming on species the underlying data has marked as weak performers in this report.

\section*{References}

\begin{enumerate}
    \item GBIF.org. Global Biodiversity Information Facility Occurrence Download (California, Asteraceae / Poaceae / Apiaceae). \textit{[placeholder: insert DOI from the GBIF download page used to seed} \texttt{data/CA\_3Species\_Raw.csv}\textit{]}.
    \item Van Horn, G. et al.\ The iNaturalist Species Classification and Detection Dataset (iNat 2021). \textit{[placeholder: full citation]}.
    \item He, K., Zhang, X., Ren, S., Sun, J. Deep Residual Learning for Image Recognition. CVPR 2016.
    \item Mac Aodha, O. et al.\ Presence-Only Geographical Priors for Fine-Grained Image Classification. ICCV 2019. \textit{[placeholder: confirm citation --- the closest analogue to this project's approach]}.
    \item Paszke, A. et al.\ PyTorch: An Imperative Style, High-Performance Deep Learning Library. NeurIPS 2019.
    \item Pedregosa, F. et al.\ Scikit-learn: Machine Learning in Python. JMLR 12 (2011).
    \item Project repository: \texttt{plant-geopriors} (this work). \textit{[placeholder: GitHub URL]}.
\end{enumerate}

\end{document}
